\chapter{Kết luận và Hướng phát triển}
\label{ch:conclusion}

Chương kết luận này tổng hợp lại toàn bộ hành trình nghiên cứu của luận văn với đề tài ``Kiến trúc nhận thức hai tầng tự động hóa phát hiện đe dọa và phản hồi sự cố sử dụng AI tác tử'' (SENTINEL). Chương này đúc rút các kết quả nghiên cứu đã được thực nghiệm chứng minh tương ứng với các câu hỏi nghiên cứu cốt lõi, tự đánh giá các giới hạn kỹ thuật và thực tiễn trong khuôn khổ một luận văn thạc sĩ, đồng thời mở ra những hướng nghiên cứu tương lai đầy hứa hẹn cho việc bảo vệ hạ tầng số hiện đại bằng AI Tác tử cục bộ được bảo vệ an toàn.

\section{Tổng hợp Kết quả Nghiên cứu Đạt được}
Tổng hợp kết quả nghiên cứu đạt được khẳng định sự giải quyết thành công ba câu hỏi nghiên cứu cốt lõi được đặt ra ở phần đầu của nghiên cứu này.

Câu hỏi thứ nhất (RQ1) liên quan đến bộ lọc ở tốc độ đường truyền phi trạng thái. Việc thiết kế và triển khai thuật toán thống kê trực tuyến Welford ở Tầng Truyền tải (Tầng 1) đã giải quyết thành công các thách thức của việc xử lý dữ liệu đo lường từ xa dung lượng lớn và nạn bội thực cảnh báo. Bằng cách duy trì đường cơ sở phiên cuộn và tính toán điểm Z-Score với độ phức tạp thời gian $O(1)$ và bộ nhớ $O(1)$, Tầng 1 đạt tỷ lệ lọc bỏ log lành tính lên tới $99.8\%$ trên bộ dữ liệu CSE-CIC-IDS2018. Sự cắt giảm đáng kể này đảm bảo rằng chỉ có $0.2\%$ các sự kiện mạng bất thường được leo thang lên tầng nhận thức tiêu tốn nhiều tài nguyên tính toán. Hơn nữa, do Tầng 1 ánh xạ đường cơ sở của các phiên hoạt động trên cơ sở từng địa chỉ IP, hệ thống bảo toàn chuỗi sự kiện thời gian của các cuộc tấn công Đe dọa Thường trực Nâng cao (APT) dạng ``chậm và ẩn nấp'' mà không làm mất đi ngữ cảnh chuỗi sự kiện cần thiết cho đối chiếu tương quan dài hạn.

Câu hỏi thứ hai (RQ2) đề cập đến phân loại LLM an toàn và giảm thiểu độ trễ. Việc tích hợp mô hình ngôn ngữ lớn lượng tử hóa cục bộ (\textit{Gemma-2-9B-IT} Q6\_K) thông qua máy trạng thái FSM có trạng thái được LangGraph điều phối giúp giải quyết đồng thời thách thức về độ trễ vận hành và lỗ hổng nhận thức. Trong khi việc nạp trực tiếp log thô vào LLM gặp phải nút thắt cổ chai do độ phức tạp bình phương của cơ chế tự chú ý (self-attention), kiến trúc SENTINEL giảm thiểu rủi ro này bằng cách chỉ định tuyến các dị biệt đã được lọc trước và sử dụng Bộ đệm Ngữ nghĩa trong bộ nhớ (dựa trên thuật toán LRU với khóa SHA-256 của các khuôn mẫu sự kiện). Bộ đệm này giải quyết các cảnh báo trùng lặp trong vài micro giây, trong khi các dị biệt chưa được lưu cache đạt độ trễ suy luận trung bình từ $4$--$6$ giây cho mỗi batch trên phần cứng máy khách. Về mặt an toàn, lớp Đóng gói Dữ liệu phân tách (tích hợp các mã token thập lục phân ngẫu nhiên được sinh bằng mật mã bảo mật) và các bộ lọc làm sạch đầu ra (Output Sanitization) chứng minh hiệu quả cao trong việc ngăn chặn các cuộc tấn công tiêm prompt từ log (S1--S4) và buôn lậu dấu phân tách, ngăn không cho tải tin độc hại chiếm quyền điều khiển ngữ cảnh lập luận của tác tử.

Câu hỏi thứ ba (RQ3) xem xét ràng buộc ngữ cảnh, bộ nhớ và kiểm chứng thống kê. Sự kết hợp của cơ chế truy xuất Dual-RAG và cơ sở dữ liệu Bộ nhớ Đe dọa SQLite cục bộ ràng buộc hiệu quả các quyết định của tác tử và giảm thiểu hiện tượng ảo giác thực tế. Sử dụng tìm kiếm vector dày đặc FAISS và tìm kiếm từ khóa thưa BM25 được tổng hợp qua thuật toán RRF ($k=60$), cỗ máy RAG truy xuất các kỹ thuật MITRE ATT\&CK và kịch bản NIST phù hợp. Sự ràng buộc này được xác thực qua quy trình đánh giá LLM-Trọng tài khác họ, trong đó các chỉ số kiểu RAGAS đạt điểm số Faithfulness là $4.00$, Answer Relevancy đạt $4.63 \pm 0.80$, Context Recall đạt $3.99 \pm 0.07$, và Context Precision đạt $3.00 \pm 0.10$ trên thang điểm từ $1$--$5$. Sự tương quan APT dài hạn được thực hiện thông qua các lược đồ Bộ nhớ Đe dọa (\texttt{ip\_reputation}, \texttt{known\_entities}, \texttt{threat\_events}, và \texttt{apt\_indicators}), nâng mức cảnh báo thành công khi một IP nguồn kích hoạt dị thường trải trên nhiều ngày hoặc nhiều giai đoạn trong chuỗi tấn công (kill-chain). Các kiểm định thống kê phi tham số củng cố vững chắc các phát hiện này bằng toán học: kiểm định McNemar cho ra p-value là $1.0$ (chứng minh Tầng 2 gia tăng năng lực giải thích mà không làm thay đổi phân loại F1 cơ sở), và kiểm định Mann-Whitney U cho ra $U=27053.5$ với $p = 2.84 \times 10^{-17}$, khẳng định sự cải thiện tốc độ từ thiết kế hai tầng là có ý nghĩa thống kê vượt trội.

\section{Đóng góp Khoa học và Thực tiễn}
Luận văn đem lại cả đóng góp khoa học và thực tiễn thiết thực cho lĩnh vực an ninh mạng.

\subsection{Đóng góp Khoa học}
Các đóng góp mới về mặt lý thuyết và khoa học của công trình này gồm ba điểm. Thứ nhất là việc hình thức hóa một kiến trúc nhận thức an ninh hai tầng giúp phân tách tầng lọc dữ liệu phi trạng thái tốc độ cao (Tầng 1) khỏi tầng lập luận có trạng thái nhận thức chuyên sâu (Tầng 2), thiết lập một mô hình kết hợp giữa toán học thống kê xác định với các mô hình ngôn ngữ lớn bất tất định. Thứ hai là việc xây dựng một mô hình cô lập mật mã cho các khuôn mẫu prompt---Đóng gói Dữ liệu Phân tách---giải quyết triệt để lỗ hổng của LLM trước các dữ liệu log không đáng tin cậy mà không cần dựa vào các bộ phân loại rào chắn mỏng manh và tốn kém tài nguyên. Thứ ba là việc đề xuất một khung đánh giá phi tham số cho các tác tử an ninh mạng cục bộ, sử dụng trọng tài LLM chéo họ và các chỉ số RAGAS để đo lường mức độ đồng bộ nhận thức, chất lượng lập luận và sự ràng buộc ngữ cảnh.

\subsection{Đóng góp Thực tiễn}
Các đóng góp về mặt kỹ thuật và ứng dụng thực tiễn của công trình này cũng gồm ba điểm. Thứ nhất là một nguyên mẫu hoàn chỉnh mã nguồn mở (SENTINEL) được viết bằng Python và đóng gói qua Docker Compose, tích hợp \texttt{llama.cpp} và \texttt{LangGraph} giúp doanh nghiệp nhanh chóng thử nghiệm thực tế. Thứ hai là việc xác thực khả năng phân loại đe dọa cục bộ biệt lập (air-gapped) trên phần cứng thương mại phổ thông---gồm CPU Intel Core i7-14700KF, 32\,GB RAM DDR5 và một GPU NVIDIA RTX 4060 Ti duy nhất---đảm bảo chủ quyền dữ liệu và tuân thủ các nguyên tắc Zero-Trust cho các tổ chức có tài nguyên hạn chế. Thứ ba là một bảng điều khiển quản trị (Streamlit UI) kết nối tự động hóa với sự giám sát của con người thông qua hàng đợi HITL và bảo vệ tệp tin nhật ký an ninh trước các hành vi can thiệp sửa đổi bằng kỹ thuật liên kết chuỗi HMAC.

\section{Các Giới hạn Hiện tại}
Mặc dù hệ thống đã chứng minh tính ưu việt rõ rệt, nghiên cứu vẫn thẳng thắn thừa nhận một số giới hạn kỹ thuật phân định phạm vi hoạt động. Về giới hạn dữ liệu và định dạng, thực nghiệm mới chỉ được đánh giá chủ yếu trên log luồng mạng (Network Flow) và siêu dữ liệu HTTP; hệ thống thiếu các bộ phân tích (parser) cấu hình sẵn cho các nguồn log phi cấu trúc đa dạng trong doanh nghiệp---như Windows Event Logs, Linux Syslogs hay CloudTrail---vốn đòi hỏi thêm các thiết lập trích xuất khuôn mẫu log Drain riêng biệt. Về vấn đề bão hòa hàng đợi dị thường, mặc dù bộ lọc Welford lọc bỏ $99.8\%$ lưu lượng lành tính, một cuộc tấn công mạng quy mô lớn như quét cổng tốc độ cao hoặc tấn công từ chối dịch vụ, khi tạo ra lượng dị thường tăng vọt đột biến, sẽ làm bão hòa hàng đợi leo thang và dẫn đến giới hạn dung lượng VRAM cùng độ trễ xếp hàng nếu tài nguyên GPU không được mở rộng kịp thời. Một sự đánh đổi khác về tự động hóa nảy sinh do, để ngăn nguy cơ tự từ chối dịch vụ (Self-DoS) ngoài ý muốn khi LLM có thể chặn nhầm các tài sản nội bộ quan trọng (như Active Directory hoặc máy chủ cơ sở dữ liệu), kiến trúc phải dựa vào cơ chế duyệt của con người (HITL) để thực thi chặn tường lửa; điều này bảo đảm an toàn nhưng làm hạn chế tốc độ phản phó hoàn toàn tự động theo thời gian thực. Về tấn công đầu độc ngữ nghĩa trực tuyến, mặc dù việc kiểm tra mã băm checksum của cơ sở dữ liệu vector ngăn chặn hành vi can thiệp ngoại tuyến vào tệp tin, hệ thống vẫn có thể bị tổn thương nếu kẻ tấn công có quyền ghi vào thư mục tri thức thông qua một kênh được cấp phép để tiêm vào các tài liệu hướng dẫn (playbook) sai lệch. Cuối cùng, bản thân quá trình đánh giá bị giới hạn về phạm vi và năng lực thống kê, nên được hiểu là mang tính chỉ báo chứ không phải toàn diện: năng lực zero-day được kiểm chứng qua một proxy thống kê có kiểm soát---các luồng lành tính bị đẩy tới các độ lệch Welford phân cấp, với ranh giới phát hiện được xác định ở $\approx 4\sigma$---chứ không phải các mẫu mã độc thực tế; các nghiên cứu APT, đối kháng và chất lượng suy luận dựa trên các mẫu khiêm tốn (ba chiến dịch, báo kèm khoảng tin cậy Wilson $95\%$ là $[0.44, 1.00]$ và một đối chứng âm không-báo-nhầm, cùng 120 tải trọng và một tập con phân tầng 300 sự kiện); nghiên cứu loại trừ đầy đủ sáu cấu hình (A--F) được đánh giá trên cả một tập con thiên-tấn-công lẫn một tập con cân bằng $150/150$---tập sau cô lập đóng góp precision và độ trễ của cổng Welford so với baseline LLM thuần---song cả hai đều dựa trên một mô hình lượng tử hóa duy nhất trên một máy chủ duy nhất, nên các con số tuyệt đối có thể thay đổi với các mô hình, phần cứng hoặc thành phần lưu lượng khác.

\section{Hướng Phát triển Tương lai}
Từ các giới hạn trên, một số hướng phát triển tương lai được định hình để tiếp nối nghiên cứu này. Hướng thứ nhất là phát triển hệ thống đa tác tử (Multi-Agent Systems - MAS), phân rã tác tử LangGraph đơn khối thành một mạng lưới các tác tử chuyên biệt hợp tác với nhau---gồm các tiểu tác tử như Tác tử Trích xuất IOC, Tác tử Đối chiếu Playbook, Tác tử Soạn thảo Quy tắc Tường lửa và Tác tử Lập Báo cáo Pháp y---nhằm tối ưu hóa chất lượng lập luận. Hướng thứ hai là học liên hợp bảo vệ quyền riêng tư (Federated Learning), mở rộng kiến trúc sang mô hình phi tập trung, nơi nhiều nút SENTINEL phân tán có thể chia sẻ các trọng số mô hình hoặc dị biệt uy tín đã học được cục bộ mà không cần truyền tải các log thô nhạy cảm vi phạm quyền riêng tư. Hướng thứ ba là học tăng cường từ phản hồi của chuyên gia, tích hợp một vòng lặp học tăng cường cục bộ, trong đó tác tử cập nhật các tham số sinh quy tắc tường lửa dựa trên các hành động chấp thuận hoặc từ chối rõ ràng của chuyên viên SOC trong hàng đợi Streamlit HITL. Hướng thứ tư là tự thích ứng đối kháng (Adversarial Self-Hardening), triển khai một tác tử giả lập tấn công (Red Team Agent) chuyên biệt trong mạng container để tự động thăm dò SENTINEL bằng các cuộc tấn công tiêm prompt và vượt rào, qua đó giúp hệ thống tự động kiểm thử và củng cố các token rào chắn của chính mình. Hướng thứ năm là định danh chiến dịch dựa trên TTP (TTP-based Campaign Attribution). Hệ thống hiện tại phát hiện \emph{sự tồn tại} của một chiến dịch APT đa-ngày, đa-giai-đoạn theo cơ chế nổi lên (emergent), phi giám sát---thông qua đối chiếu tương quan đa-ngày trong Bộ nhớ Đe dọa---nhưng chủ ý KHÔNG gọi tên nhóm tấn công (intrusion set) cụ thể đứng sau. Một mở rộng tự nhiên là bổ sung một mô-đun sinh giả thuyết định danh nằm trên các định danh kỹ thuật có cấu trúc mà Lớp ánh xạ ATT\&CK đã sản sinh: mỗi chuỗi kỹ thuật quan sát được theo từng nguồn sẽ được đối chiếu với hồ sơ kỹ thuật đã công bố của các Nhóm MITRE ATT\&CK (ATT\&CK Groups) bằng một độ đo tương đồng tập có trọng số---ví dụ Jaccard hoặc chồng lấp trọng số kiểu TF-IDF nhằm nâng trọng số cho các kỹ thuật hiếm---cho ra một danh sách các chiến dịch ứng viên được xếp hạng, mỗi ứng viên kèm một độ tin cậy tương đồng. Vì mở rộng này tái dùng chính đầu ra kỹ thuật có cấu trúc và chuỗi kill-chain theo từng nguồn vốn đã được duy trì trong Bộ nhớ Đe dọa, nó chủ yếu chỉ đòi hỏi việc nạp hồ sơ các Nhóm ATT\&CK vào cơ sở tri thức. Tuy nhiên, điều cốt yếu là đầu ra này phải được trình bày nghiêm ngặt như một GIẢ THUYẾT được xếp hạng để điều tra viên xem xét, chứ không phải một phán quyết dứt khoát: định danh trong thực tế phụ thuộc vào tình báo vượt xa các thủ đoạn dùng chung---hạ tầng, thời điểm, dấu vết ngôn ngữ, bối cảnh địa chính trị, và khả năng luôn hiện hữu của việc cố ý giả mạo cờ hiệu (false-flag)---và do các kỹ thuật thường được dùng chung giữa nhiều nhóm, một phép khớp thuần dựa trên TTP chỉ có thể là bằng chứng bổ trợ, không bao giờ là bằng chứng khẳng định về tác giả.

\subsection{Lộ trình Mở rộng Quy mô theo Chiều ngang}
Hướng thứ sáu mang tính xuyên suốt là khả năng mở rộng quy mô theo chiều ngang (Horizontal Scalability). Nguyên mẫu hiện tại được chủ ý triển khai trên một máy chủ đơn qua Docker Compose, nên thừa hưởng các nút thắt cổ chai của kiến trúc đơn nút mà hai giới hạn đã nêu ở trên---bão hòa hàng đợi dị thường và năng lực tính toán trên một GPU duy nhất---làm bộc lộ rõ. Một bước tiến hóa lên môi trường sản xuất sẽ đi theo cẩm nang kinh điển của hệ phân tán trên năm trục phối hợp, được lựa chọn sao cho giao ước nhận thức hai tầng cốt lõi vẫn giữ nguyên.

Trục thứ nhất là \textit{phân mảnh (sharding) tầng lọc Tier-1 có trạng thái}. Tier-1 hiện lưu đường cơ sở Welford theo từng địa chỉ IP nguồn trong bộ nhớ của một tiến trình duy nhất, khiến nó trở thành một dịch vụ có trạng thái chưa thể mở rộng theo chiều ngang. Do trạng thái theo thời gian của mỗi IP nguồn độc lập với mọi IP khác, bộ lọc có thể được phân mảnh trên một cụm nhiều nút bằng cách áp dụng \emph{băm nhất quán (consistent hashing)} lên khóa IP nguồn: mỗi IP được định tuyến một cách tất định tới đúng nút đang nắm giữ đường cơ sở của nó, nhờ đó bảo toàn nguyên vẹn đường cơ sở theo từng IP lẫn chuỗi sự kiện APT ``chậm và ẩn nấp''; năng lực lọc tăng tuyến tính theo số nút; và việc thêm hay bớt một nút chỉ tái cân bằng $k/n$ khóa thay vì toàn bộ không gian khóa. Đây chính là bước biến thành phần có-trạng-thái duy nhất thành một tầng mở rộng được theo chiều ngang.

Trục thứ hai là \textit{nhật ký thu nhận được phân mảnh và nhân bản}. Bộ môi giới Redis Streams đơn dùng trong nguyên mẫu là một điểm lỗi đơn (single point of failure); nó sẽ được thay bằng một nhật ký cam kết (commit log) được phân mảnh và nhân bản (Redis Cluster hoặc Apache Kafka), với khóa phân mảnh được căn chỉnh trùng khóa sharding của Tier-1, để mọi sự kiện của một nguồn nhất quán đi tới đúng nút lọc đang sở hữu đường cơ sở của nguồn đó. Trục thứ ba là \textit{tiêu thụ bất biến (idempotent)}: vì nhóm tiêu thụ (consumer group) hiện đã bảo đảm phân phát ít-nhất-một-lần (at-least-once), một lớp bất biến khóa theo một định danh sự kiện tất định sẽ được bổ sung, để việc phân phát lại sau sự cố không thể thực thi trùng lặp một hành động chặn.

Trục thứ tư là \textit{cụm nhân công nhận thức co giãn (elastic worker pool)}. Thông lượng Tier-2 hiện bị chặn trần bởi một GPU duy nhất; nó sẽ được mở rộng thành một cụm các nhân công suy luận LLM phi trạng thái, cùng rút việc từ hàng đợi leo thang, với số lượng nhân công được tự động co giãn theo độ sâu hàng đợi quan sát được. Điều này trực tiếp giảm nhẹ giới hạn bão hòa hàng đợi---một đợt bùng nổ lưu lượng làm hàng đợi dài ra, và chính điều đó lại kích hoạt cấp phát thêm năng lực suy luận---trong khi Tier-1 tất định vẫn tiếp tục bảo vệ mạng một cách độc lập trong lúc suy luận bị dồn ứ. Trục thứ năm là \textit{tầng lưu trữ bền vững được nhân bản}. Bộ nhớ Đe dọa và kho audit SQLite sẽ được di trú sang một tầng dữ liệu chịu lỗi với mô hình nhất quán chọn theo từng loại tải: một kho khóa-giá-trị dựa trên đồng thuận số đông (thiết kế kiểu Dynamo hoặc Cassandra với ngưỡng $W+R>N$ tinh chỉnh được) phù hợp với tải uy tín IP và đối chiếu APT vốn dung thứ nhất quán cuối cùng (eventual consistency), trong khi chuỗi audit HMAC chống giả mạo và kho luật do con người phê duyệt vẫn giữ nhất quán mạnh trên một cơ sở dữ liệu quan hệ được nhân bản.

Đặt một bộ cân bằng tải trước tầng bảng điều khiển phi trạng thái và một bộ giới hạn tốc độ (rate limiter) phía thu nhận đứng trước bộ môi giới để hấp thụ các trận lũ lưu lượng sẽ hoàn thiện bức tranh. Gộp lại, năm trục này vạch ra một lộ trình di trú cụ thể từ một nguyên mẫu nghiên cứu đơn-máy đã được kiểm chứng sang một triển khai chịu lỗi, mở rộng được theo chiều ngang---giải quyết đúng các giới hạn bão hòa và một-GPU đã nêu ở trên---mà không làm thay đổi kiến trúc nhận thức hai tầng mà luận văn này đã xác thực.

\vspace{1cm}
\begin{flushright}
    \textit{Khép lại một nỗ lực khiêm tốn, nhưng mở ra những khát vọng sâu sắc trong hành trình bảo vệ không gian số bằng sức mạnh của AI Tác tử.}
\end{flushright}
